% Generated draft: requires mathematical and source review.
\documentclass[11pt]{article}
\usepackage[margin=25mm]{geometry}
\usepackage{fontspec}
\setmainfont{Latin Modern Roman}
\usepackage{amsmath,amssymb,mathtools,unicode-math}
\setmathfont{Latin Modern Math}
\usepackage{enumitem,needspace,xurl,hyperref}
\setlength{\parindent}{0pt}
\setlength{\parskip}{5pt}
\setlength{\emergencystretch}{3em}
\newcommand{\N}{\mathbb{N}}
\newcommand{\Z}{\mathbb{Z}}
\newcommand{\Q}{\mathbb{Q}}
\newcommand{\R}{\mathbb{R}}
\newcommand{\C}{\mathbb{C}}
\newcommand{\F}{\mathbb{F}}
\newcommand{\A}{\mathbb{A}}
\newcommand{\PP}{\mathbb{P}}
\newcommand{\E}{\mathbb{E}}
\newcommand{\eps}{\varepsilon}
\newcommand{\dd}{\,\mathrm{d}}
\begin{document}
\section*{UnsolvedMath ID: 30005142}
\textbf{Source problem code:} catalogue entry
\subsection{Definitions and mathematical statement}
Let \(M\) denote the manifold with boundary referred to in the source, and let \(n=\dim M\). Neither \(M\), its dimension, nor its geometric hypotheses are specified in the supplied excerpt. For a Riemannian metric \(g\) on \(M\), define
\[
V(g):=\operatorname{Vol}_{n}(M,g),\qquad
A(g):=\operatorname{Vol}_{n-1}(\partial M,g|_{T\partial M}),\qquad
R(g):=\frac{A(g)}{V(g)},
\]
whenever these volumes are finite and \(V(g)>0\). A positive prescribed ratio also requires \(A(g)>0\).

Write \(\mathcal{Y}_{\mathrm{source}}(M)\) for the class of metrics intended by the source's term ``Yamabe metrics.'' This is a placeholder for missing source context: the supplied record does not specify a background conformal class, a normalization, the relevant variational characterization, or the scalar-curvature and boundary conditions defining this class.

The extracted question is whether the analysis cited in the report as reference [9] can be refined to produce a family
\[
(g_T)_{T\in(0,\infty)},\qquad g_T\in\mathcal{Y}_{\mathrm{source}}(M),
\]
such that
\[
\forall T\in(0,\infty),\qquad
\frac{\operatorname{Vol}_{n-1}(\partial M,g_T|_{T\partial M})}
{\operatorname{Vol}_{n}(M,g_T)}=T.
\]
The intended regularity of the dependence on \(T\) is unspecified; the displayed indexing alone imposes no continuity or smoothness. The original excerpt writes \(g_u\) in the volume quotient while indexing the proposed family by \(T\); the supplied cleaned statement uses \(g_T\) throughout. The relationship between \(u\) and \(T\) is not given. Because the quotient depends on the scale of the metric, any intended normalization or restriction on rescaling is material. Thus the displayed question remains schematic until the omitted hypotheses and the meaning of ``Yamabe metrics'' are recovered.
\subsection{Short English statement}
The question asks whether the analysis cited as [9] can produce Yamabe metrics whose boundary-volume to interior-volume ratio takes every value \(T>0\). The supplied excerpt leaves the manifold hypotheses, the precise class of Yamabe metrics, metric normalization, and parameter regularity unresolved. It therefore does not determine a complete mathematical formulation.
\subsection{Sources}
\begin{itemize}
\item Supplied local UnsolvedMath registry record, permanent ID \(30005142\), title ``Yamabe Metrics with Prescribed Boundary Volume Ratio,'' source problem code \texttt{catalogue entry}, in the collection ``Oberwolfach Reports Open Problems.'' No individual proposer is named. The record gives proposed year 2022.
\item Supplied report citation: \emph{Geometrie (2023)}, Oberwolfach Reports, DOI 10.4171/owr/2022/28, pp. 1551--1601. \url{https://doi.org/10.4171/owr/2022/28}. The citation's 2023 label and the record's 2022 proposed year are retained without reconciliation.
\item Supplied literature source, ``OWR report'': \url{https://ems.press/content/serial-article-files/46965}. The excerpt refers to analysis in [9], but the bibliographic entry for [9] is absent from the supplied record.
\item Supplied literature source: Escobar, \emph{The Yamabe problem on manifolds with boundary}. \url{https://doi.org/10.4310/JDG/1214447805}.
\item Supplied literature source: Brendle--Chen, ``boundary Yamabe existence.'' \url{https://arxiv.org/abs/0908.4327}.
\item Literature assessment dated exactly \texttt{2026-08-22}: the supplied curation labels the problem open and describes the prescribed all-positive volume-ratio family as an unproved refinement question. It states that standard boundary Yamabe existence and local uniqueness results do not establish this global one-parameter family. This is supplied dated curation derived from the OWR record, not an independently verified assessment; its time-sensitive status and cited claims require independent verification.
\end{itemize}
\end{document}
